\section{Methodology}
\label{sec:methodology}

\subsection{Data Collection and Curation}
Our Harappan database integrates two primary sources:
\begin{itemize}
  \item Mahadevan's concordance (417 signs)
  \item Wells' extended inventory (694 signs)
\end{itemize}

Each sign is annotated with frequency data, positional statistics,
and proposed phonetic/semantic values from the literature.

\subsection{Phoneme Mapping Approach}
Universal phoneme mappings are constructed by:
\begin{enumerate}
  \item Extracting phoneme inventories from language documentation
  \item Aligning features using IPA as an intermediary
  \item Building cross-linguistic correspondence tables
\end{enumerate}

\subsection{Symbolic Layer Construction}
The symbolic/numerological layer applies:
\begin{itemize}
  \item Pythagorean digit reduction
  \item Hebrew gematria principles
  \item Vedic numerological correspondences
\end{itemize}

\subsection{Pipeline Implementation}
The processing pipeline is implemented in Python, with:
\begin{itemize}
  \item Modular functions for each stage
  \item JSON/CSV data interchange formats
  \item REST API endpoints for integration
\end{itemize}

\subsection{Evaluation Criteria}
Given the exploratory nature of this research, evaluation focuses on:
\begin{itemize}
  \item Internal consistency of mappings
  \item Reproducibility of results
  \item Utility for downstream research tasks
\end{itemize}
